\section{Method}
\label{sec:method}

\subsection{Problem formulation and overview}

Our goal is to improve the expected physical utility of generated accelerator
RTL at a fixed draw budget. For specification $d$, policy $\pi_\theta$ samples completion
$m$. Let $g_d(m)$ be the binary executable-oracle verdict and $F(m)$ the
post-route timing-derived frequency, defined as zero if implementation fails.
The objective over a design distribution $\mathcal{D}$ is
\begin{equation}
\max_\theta J(\theta),\qquad
J(\theta)=\operatorname{E}_{d\sim\mathcal{D}}
\operatorname{E}_{m\sim\pi_\theta(\cdot\mid d)}[g_d(m)F(m)].
\label{eq:problem}
\end{equation}
Because failures contribute zero, this objective favors policies that assign
more probability to fast implementations that pass the oracle.
Equation~(\ref{eq:equal-sample}) gives the finite-sample estimate used in the
evaluation. Implementing every training sample would be too expensive, so a
learned surrogate supplies the training reward and measured post-route timing
determines the reported result. Fig.~\ref{fig:training-boundary} places these
steps in context. Domain SFT provides both the initial policy and a frozen
reference; fresh completions pass through the executable gate and receive a
physical or control reward. Group-relative updates then change the policy's
sampling distribution. The final comparison evaluates frozen checkpoints at
equal draw budgets using the oracle and Vivado.


Each accelerator family has a deterministic software reference. The oracle
compiles generated modules and compares simulated outputs on frozen stimuli.
To accommodate different pipeline depths, it searches a bounded latency window and
accepts the best aligned match fraction at or above
\revisionclaim{OracleThreshold}.  The search allows zero through
\revisionclaim{OracleMaxLatency} cycles of latency and discards
\revisionclaim{OracleWarmup} warm-up outputs. A sufficiently long trace can pass
this near-unity threshold despite a mismatch. We use ``correct'' below to mean
oracle-passing, not formally equivalent. The same oracle filters supervised examples, gates
policy rewards, and evaluates held-out samples, with held-out stimulus streams
reserved for evaluation. Let $C_d(m)$ denote the binary oracle verdict.
Every reward arm obeys
\begin{equation}
g_d(m)=
\begin{cases}
\one, & C_d(m)\text{ passes},\\
\zero, & C_d(m)\text{ fails}.
\end{cases}
\label{eq:gate}
\end{equation}
The gate rejects behavior changes detected by its sampled tests.  It neither
proves equivalence on all inputs nor establishes that a proxy ranks admissible
modules properly; physical implementation addresses the latter question.

\subsection{Physical reward and matched controls}

The physical reward ranks oracle-passing completions. It first applies a frozen
lexical canonicalizer $\kappa$, then predicts log-frequency from canonical structural
features $\phi$:
\begin{equation}
R_d^{\mathrm{RF}}(m)=g_d(m)\,
\exp\!\left[\operatorname{clip}\!\left(
\widehat{f}_{\mathrm{RF}}(\phi(\kappa(m))),a,b\right)\right].
\label{eq:rf-reward}
\end{equation}
Here $\widehat{f}_{\mathrm{RF}}$ predicts log-frequency, $a,b$ are the frozen
log-frequency clamp bounds, and exponentiation returns the reward in MHz.
The RF contains \studyclaim{StudyTwoRewardTreeCount} trees and uses
\studyclaim{StudyTwoRewardFeatureCount} features: multiplier count, maximum
per-statement multiplier count, positive-edge count, nonblocking and blocking
assignment counts, multiply--accumulate structure, register and wire counts,
ternary and comparison counts, statement count, a structural pipeline ratio,
and loop count, maximum bound, and summed bounds.  Descriptive impurity
importance ranks multiplier count, multiply--accumulate structure, pipeline
ratio, register count, and nonblocking assignments highest; this ranking
describes the fitted predictor and is not a causal circuit attribution.  It is fitted to
\studyclaim{StudyTwoRewardTrainingRows} measured RTL rows that exclude every
held-out and sealed design.  Layout-dependent line counts and name-dependent
features are absent.  The canonicalizer is lexical and layout-invariant under
its frozen mutation suite; it is not a canonical hardware representation.
Canonicalization or canonical-compilation failure also sets the reward to zero
and is logged.

An earlier reward failure motivated the choice of canonical features and the
RF predictor, as documented in the Supplementary Material. We fixed these
features and hyperparameters before opening the sealed efficacy split, which
allowed the repair hypothesis to be tested prospectively on that split.
The hypothesis itself was informed by the earlier outcomes. The RF provides
reward variation in \studyclaim{StudyTwoRewardResolution}\% of eligible
training groups. The supplement also records the reward-path validation and the
changes to its scope before sealed evaluation.

Two controls assess the role of the reward. The original-MLP arm restores the
previous textual/structural predictor while keeping the SFT initialization,
prompts, oracle, optimizer, SFT reference, update target, and evaluation fixed.
The correctness-only arm uses $R_d^{\mathrm{C}}(m)=g_d(m)$ and therefore tests
whether correctness pressure suffices in the realized training run.  Because
canonicalization, features, and predictor class change jointly, RF versus MLP
is not a component-wise ablation.  The RF arm
has \studyclaim{StudyTwoRfTrainingSeeds} independent training seeds, the MLP
arm has \studyclaim{StudyTwoMlpTrainingSeeds}, and the correctness-only arm has
\studyclaim{StudyTwoCorrectnessTrainingSeeds} preregistered full-budget seed.

\subsection{Supervised initialization and policy training}

We adapt RTLCoder~\cite{liu2024rtlcoder} on an oracle-verified accelerator
corpus using completion-masked LoRA~\cite{hu2022lora}.  Held-out designs are
excluded from the corpus, reward-training rows, and policy prompts.  The
supervised adapter supplies both the policy initialization and the frozen
reference, allowing optimization to remain anchored to the domain competence
established by supervision. Using the raw base model as the reference would
instead regularize toward a model that had not received this domain adaptation.

Each training group contains \studyclaim{StudyTwoGroupSize} completions sampled
at temperature \studyclaim{StudyTwoTrainingTemperature}, with at most
\studyclaim{StudyTwoTrainingMaxTokens} completion tokens.  If $R_i$ is the
gated reward of completion $m_i$, its within-group advantage is
\begin{equation}
A_i=\frac{R_i-\overline{R}}{s_R+\epsilon}.
\label{eq:advantage}
\end{equation}
We sample each attempted group from the current policy and use its rewards
directly in the update. Groups without reward variation provide no ranking
information and do not advance the optimizer-update counter. For an informative group, the
implemented objective uses the mean completion-token log probability
$\ell_\theta(m\mid d)=|m|^{-\one}\sum_t
\log\pi_\theta(m_t\mid d,m_{<t})$:
\begin{equation}
\mathcal{L}_i=-A_i\ell_\theta(m_i\mid d)
+\beta\left(\exp(\Delta_i)-\Delta_i-\one\right),
\label{eq:policy-loss}
\end{equation}
where $\Delta_i=\ell_{\mathrm{SFT}}(m_i\mid d)-
\ell_\theta(m_i\mid d)$ and the reference is detached.  Losses are averaged
within the group.  This length-normalized update is inspired by group-relative
policy optimization~\cite{shao2024deepseekmath}; it is not the standard clipped
GRPO objective.  Exponentiating a difference of mean token log probabilities
does not yield the usual sequence-probability ratio, so the penalty is a
reference regularizer rather than an unbiased KL estimate.
Every arm uses AdamW with learning rate
\studyclaim{StudyTwoLearningRate}, reference-penalty coefficient
\studyclaim{StudyTwoKlCoefficient}, and a target of
\studyclaim{StudyTwoTrainingUpdates} non-flat updates. Checkpoints follow the
training update counter rather than attempted groups; the frozen midpoint is
update \studyclaim{StudyTwoMidpointUpdate}.

\subsection{Physical evaluation and distribution analysis}

Evaluation implements each distinct oracle-passing candidate once and restores
its original sample multiplicity.  For policy $p$ and design $d$,
let $\mathcal{M}_{p,d}$ be the distinct correct candidates, $c_m$ the sampled
multiplicity, $n_d$ the complete draw budget, and $F(m)$ the post-route
timing-derived frequency, defined as zero when implementation fails.  The
primary per-design endpoint is
\begin{equation}
\overline{F}_{p,d}=\frac{\sum_{m\in\mathcal{M}_{p,d}}c_mF(m)}{n_d}.
\label{eq:equal-sample}
\end{equation}
Incorrect samples are absent from $\mathcal{M}_{p,d}$ but remain in $n_d$ and
therefore contribute zero.  Each arm receives
\studyclaim{StudyTwoSamplesPerArmDesign} draws per design; a two-seed arm uses
\studyclaim{StudyTwoSamplesPerSeedDesign} draws from each seed.  The headline
statistic is the unweighted mean of $\overline{F}_{p,d}$ across designs, so the
design rather than the number of distinct candidates is the aggregation unit.
Each candidate is implemented once under the same frozen clock request; $F(m)$
is therefore derived from that request and worst negative slack rather than
from candidate-specific repeated timing closure.  This symmetry supports the
paired comparison but bounds absolute-frequency claims to the frozen flow,
device, and request.

The endpoint decomposes as
\begin{equation}
\overline{F}_{p,d}=q_{p,d}\,\mu_{p,d},
\label{eq:decomposition}
\end{equation}
where $q_{p,d}$ is empirical correctness and $\mu_{p,d}$ is multiplicity-
weighted frequency conditional on an oracle-passing draw, retaining zero for
implementation failure.  We set $\mu_{p,d}=\zero$ when no draw passes.  A policy
that rarely emits very fast correct RTL can have high $\mu_{p,d}$ and low
$\overline{F}_{p,d}$; reporting only the former would change the estimand.

For paired uncertainty, optimized-minus-SFT differences are computed per
design and resampled within frozen family-by-regime strata.  The primary plan
uses \studyclaim{StudyTwoBootstrapReplicates} resamples and a
\studyclaim{StudyTwoBootstrapConfidence}\% percentile interval.  The direct
RF-minus-MLP analysis reuses the identical resample plans but is labeled post-
primary.  Candidate draws are not treated as independent experimental units,
and these design-level intervals do not estimate variation from SFT training,
reward-data construction, or model-family choice.



Post-primary diagnostics examine how probability is distributed over passing
candidate strings and compare one RF draw with a perfect selector over repeated
SFT draws. This selector uses the measured finite SFT support, retains failures
as zero-valued outcomes, and assumes that physical measurements are available
before selection. It therefore describes an inference workload on the observed
support rather than a complete training-and-deployment cost comparison. The
Supplementary Material gives the exact calculation, implementation counts, and
artifact checks; training cost is excluded from these workloads.
